\documentclass[11pt,a4paper]{article}
\usepackage{amsmath,amssymb,bm}
\usepackage{geometry}
\usepackage{hyperref}
\usepackage{parskip}
\usepackage{microtype}
\usepackage{listings}
\usepackage{xcolor}
\usepackage{booktabs}
\usepackage{caption}
\usepackage{enumitem}

\geometry{margin=2.5cm}
\setlength{\parindent}{0pt}

\hypersetup{
  colorlinks=true,
  linkcolor=blue!70!black,
  urlcolor=blue!70!black,
}

\lstset{
  basicstyle=\ttfamily\small,
  keywordstyle=\color{blue},
  commentstyle=\color{green!60!black},
  breaklines=true,
  frame=single,
  language=Python,
}

\title{\textbf{VBD Sparse Articulation Solve} \\
  \large A technical explainer for the \texttt{horde/vbd-sparse-articulation-main} branch}
\author{Newton Physics Engine}
\date{June 2026}

\begin{document}
\maketitle
\tableofcontents
\newpage

% =====================================================================
\section{Background: simulating articulated robots}
% =====================================================================

\subsection{What is an articulation?}

In Newton (and robotics simulation in general) an \emph{articulation} is a
tree (or graph) of rigid bodies connected by \emph{joints}.  A humanoid robot,
for example, is one articulation: the pelvis is the root body, and every limb
segment is a child body attached to its parent by a joint (revolute, fixed,
prismatic, etc.).

Each rigid body has 6 degrees of freedom (DOFs): 3 translational and 3
rotational.  Its pose at time $t$ is a transform $\mathbf{q} = (\mathbf{p},
\mathbf{r}) \in \mathbb{R}^3 \times \mathrm{SO}(3)$.

A joint constrains the relative motion between a parent body $a$ and a child
body $b$.  Typical joint types used in Newton are:

\begin{center}
\begin{tabular}{ll}
\toprule
Joint & Constrained DOFs \\
\midrule
Fixed      & all 6 (no relative motion) \\
Revolute   & 5 (free rotation about one axis) \\
Prismatic  & 5 (free translation along one axis) \\
Ball       & 3 (free rotation about all axes) \\
D6         & user-defined mix of linear/angular DOFs \\
Cable      & slack-only: pulls when taut, free otherwise \\
\bottomrule
\end{tabular}
\end{center}

Joints are enforced via penalty (spring-like) forces.  A high stiffness makes
the joint nearly rigid; a low stiffness allows compliance.

% =====================================================================
\subsection{The VBD / AVBD rigid-body solver}

Newton's rigid-body solver uses \emph{Augmented Vertex Block Descent}
(AVBD), an iterative position-based method.  At a high level:

\begin{enumerate}[nosep]
  \item Integrate forward to a tentative ``inertia'' pose $\mathbf{q}^\star$
        (explicit Euler from the previous velocity and gravity).
  \item Iteratively correct $\mathbf{q}$ to satisfy joint and contact
        constraints via Newton steps.
  \item Extract new velocities from $(\mathbf{q} - \mathbf{q}_\text{prev})/\Delta t$.
\end{enumerate}

At each iteration the solver minimises a total energy:
\begin{equation}
  E(\mathbf{q}) = \underbrace{E_\text{inertia}(\mathbf{q})}_{\text{stay near
  inertia pose}} + \underbrace{\sum_j E_j(\mathbf{q})}_{\text{joint penalties}}
  + \underbrace{\sum_c E_c(\mathbf{q})}_{\text{contact penalties}}
  \label{eq:energy}
\end{equation}

Because $E$ is nonlinear, minimising it directly is expensive.  Newton's
method instead approximates $E$ near the current pose $\mathbf{q}$ as a
quadratic — a bowl-shaped surface — and jumps to the bottom of that bowl in
one step.  The approximation is a Taylor expansion in a small pose update
$\boldsymbol{\delta} \in \mathbb{R}^6$ (3 translational + 3 rotational):
\begin{equation}
  E(\mathbf{q} \oplus \boldsymbol{\delta})
  \approx E(\mathbf{q})
  + \underbrace{\frac{\partial E}{\partial \mathbf{q}}}_{\mathbf{g}^\top}
    \boldsymbol{\delta}
  + \frac{1}{2}\,\boldsymbol{\delta}^\top
    \underbrace{\frac{\partial^2 E}{\partial \mathbf{q}^2}}_{\mathbf{H}}
    \boldsymbol{\delta}.
  \label{eq:taylor}
\end{equation}
Here $\mathbf{g} \in \mathbb{R}^6$ is the gradient (net force/torque pulling
the body away from the minimum) and $\mathbf{H} \in \mathbb{R}^{6\times 6}$
is the Hessian, which encodes the curvature — intuitively, the
``effective stiffness'' in every direction.  The bottom of the bowl is where
the gradient of \eqref{eq:taylor} with respect to $\boldsymbol{\delta}$
vanishes, i.e.\ $\mathbf{g} + \mathbf{H}\boldsymbol{\delta} = \mathbf{0}$,
which gives the \emph{Newton step}:
\begin{equation}
  \mathbf{H}\,\boldsymbol{\delta} = -\mathbf{g}.
  \label{eq:newton_step}
\end{equation}
This is just a linear system: given the current forces and stiffnesses, solve
for the pose correction $\boldsymbol{\delta}$ that would zero the forces if
the bowl approximation were exact.  The body pose is then updated as
$\mathbf{q} \leftarrow \mathbf{q} \oplus \boldsymbol{\delta}$ and the process
repeats.

Each energy term contributes additively to both $\mathbf{g}$ and
$\mathbf{H}$.  The inertia term adds $m/\Delta t^2$ to the translational
diagonal of $\mathbf{H}$ and $\mathbf{R}\mathbf{I}_B\mathbf{R}^\top /
\Delta t^2$ to the rotational diagonal ($m$ = mass, $\mathbf{I}_B$ = inertia
tensor in body frame, $\mathbf{R}$ = current rotation matrix).  Each penalty
constraint $E_j = \tfrac{k}{2}\|\mathbf{C}\|^2$ contributes
$k\,\mathbf{J}^\top\mathbf{J}$ to $\mathbf{H}$, where $\mathbf{J} =
\partial \mathbf{C}/\partial \mathbf{q}$ is the constraint Jacobian
(a $3\times 6$ matrix mapping pose velocity to constraint velocity).

In the \textbf{local} (per-body diagonal) solve, each body $i$ assembles
and solves its own $6\times 6$ system:
\begin{equation}
  \mathbf{H}_i \,\boldsymbol{\delta}_i = -\mathbf{g}_i,
  \label{eq:newton_local}
\end{equation}
treating every other body as fixed.  It is fast — one $6\times 6$ solve per
body — but ignores cross-body coupling from joints.

% =====================================================================
\subsection{Why the local solve is slow for stiff joints}
\label{sec:problem}

Consider two bodies connected by a stiff fixed joint.  The joint energy
penalises any relative displacement $\mathbf{C}$ between the anchor points:
\begin{equation}
  E_\text{joint} = \tfrac{k}{2} \|\mathbf{C}\|^2.
\end{equation}

When $k$ is large (e.g.\ $10^5$) the energy landscape is very ``stiff'' in
the direction that separates the anchors.  The local solve moves body $a$
toward body $b$, but does not simultaneously account for how body $b$ should
move in response.  The result is poor convergence: many VBD iterations are
needed to reach the constraint manifold, especially for chains and loops.

Benchmarks on synthetic 128-body fixed chains show that the local solve
barely reduces residuals in 8 iterations, whereas the coupled solve achieves
$\sim 2\times$ residual reduction.  For high stiffness-ratio cases the
improvement is up to $117\times$.

% =====================================================================
\section{The sparse articulation solve: key idea}
% =====================================================================

The key insight is: instead of solving each body's $6\times 6$ system
independently, solve all bodies in an articulation \emph{simultaneously}.
Each body $i$ now has a $6$-vector of unknowns $\boldsymbol{\delta}_i$
(the pose correction), and all of these are coupled into one linear system:
\begin{equation}
  \underbrace{\begin{pmatrix}
    \mathbf{H}_{11} & \mathbf{H}_{12} & \cdots & \mathbf{H}_{1n} \\
    \mathbf{H}_{21} & \mathbf{H}_{22} & \cdots & \mathbf{H}_{2n} \\
    \vdots          &                 & \ddots & \vdots          \\
    \mathbf{H}_{n1} & \cdots          & \cdots & \mathbf{H}_{nn}
  \end{pmatrix}}_{\mathbf{H} \in \mathbb{R}^{6n \times 6n}}
  \underbrace{\begin{pmatrix} \boldsymbol{\delta}_1 \\ \vdots \\
  \boldsymbol{\delta}_n \end{pmatrix}}_{\boldsymbol{\delta}}
  =
  -\underbrace{\begin{pmatrix} \mathbf{g}_1 \\ \vdots \\
  \mathbf{g}_n \end{pmatrix}}_{\mathbf{g}}
  \label{eq:coupled}
\end{equation}

Each entry $\mathbf{H}_{ij} \in \mathbb{R}^{6\times 6}$ is a \emph{block}.
The diagonal block $\mathbf{H}_{ii}$ collects all forces acting on body $i$
alone: inertia, contacts, and the self-coupling of every joint at body $i$.
An off-diagonal block $\mathbf{H}_{ij}$ ($i \neq j$) is non-zero only when
bodies $i$ and $j$ share a joint; it encodes how moving body $j$ changes the
force on body $i$ through that joint.  If two bodies have no joint between
them, their off-diagonal block is exactly zero and need not be stored.

This is what \emph{block-sparse} means: most off-diagonal blocks are zero
(bodies in different limbs of a robot are not directly coupled), so only the
non-zero blocks need to be stored and computed.  The local solve is the
extreme case of discarding \emph{all} off-diagonal blocks.

To make this concrete, consider a 4-body chain $0 \to 1 \to 2 \to 3$ (root
to tip), with joints connecting $(0,1)$, $(1,2)$, and $(2,3)$.  Each joint
produces one off-diagonal block, giving a \emph{tridiagonal} block structure
($\blacksquare$ = non-zero block, $\cdot$ = zero):

\[
\mathbf{H}_\text{chain} =
\renewcommand\arraystretch{1.3}
\begin{pmatrix}
  \blacksquare & \blacksquare & \cdot        & \cdot        \\
  \blacksquare & \blacksquare & \blacksquare & \cdot        \\
  \cdot        & \blacksquare & \blacksquare & \blacksquare \\
  \cdot        & \cdot        & \blacksquare & \blacksquare
\end{pmatrix}
\qquad
\mathbf{H}_\text{local} =
\begin{pmatrix}
  \blacksquare & \cdot        & \cdot        & \cdot        \\
  \cdot        & \blacksquare & \cdot        & \cdot        \\
  \cdot        & \cdot        & \blacksquare & \cdot        \\
  \cdot        & \cdot        & \cdot        & \blacksquare
\end{pmatrix}
\]

The right matrix is what the local solve computes: four independent $6\times
6$ systems with no cross-body information.

Now consider a \emph{closed-loop} articulation: add a joint between body~3
and body~0 (closing the chain into a ring).  Two additional off-diagonal
blocks appear.  The system is solved via Cholesky factorisation
$\mathbf{H} = \mathbf{L}\mathbf{L}^\top$, which eliminates variables one by
one.  Eliminating body~0 couples bodies~1 and~3 together, creating a
\emph{fill-in} block at position $(3,1)$ that was zero in $\mathbf{H}$ but
becomes non-zero in $\mathbf{L}$ (marked $\circ$):

\[
\mathbf{H}_\text{loop} =
\renewcommand\arraystretch{1.3}
\begin{pmatrix}
  \blacksquare & \blacksquare & \cdot        & \blacksquare \\
  \blacksquare & \blacksquare & \blacksquare & \circ        \\
  \cdot        & \blacksquare & \blacksquare & \blacksquare \\
  \blacksquare & \circ        & \blacksquare & \blacksquare
\end{pmatrix}
\quad \text{(}\circ\text{ = fill-in in } \mathbf{L}\text{, zero in }\mathbf{H}\text{)}
\]

The implementation pre-computes exactly which blocks are non-zero — including
fill-in — \emph{once} at solver construction time, before the simulation
loop starts.  At runtime only those blocks are allocated and touched,
keeping both memory and computation minimal.

% =====================================================================
\section{Implementation: step by step}
% =====================================================================

\subsection{Step 1: build the sparse layout (offline, once at construction)}
\label{sec:layout}

\textbf{File}: \texttt{rigid\_sparse\_articulation.py}

Before the simulation loop starts, Newton pre-computes the \emph{sparsity
pattern} of the block matrix $\mathbf{H}$ — i.e.\ which $6\times 6$ blocks
are non-zero.

Because $\mathbf{H}$ is symmetric positive definite (inertia and stiffness
are always positive, so the energy bowl curves upward in every direction),
Cholesky factorisation applies: $\mathbf{H} = \mathbf{L}\mathbf{L}^\top$
where $\mathbf{L}$ is lower-triangular.  Since $\mathbf{H}$ is symmetric,
the upper triangle is the mirror image of the lower triangle, so we only
need to track and store the lower-triangular sparsity pattern.  The pattern
of $\mathbf{L}$ is computed symbolically via
\texttt{\_symbolic\_cholesky\_pattern}:

\begin{lstlisting}
def _symbolic_cholesky_pattern(body_count, edges):
    rows = [{i} for i in range(body_count)]
    for a, b in edges:
        row, col = max(a,b), min(a,b)
        rows[row].add(col)                 # add the joint edge
    for k in range(body_count):
        later = [i for i in range(k+1, body_count) if k in rows[i]]
        for i_index, i in enumerate(later):
            for j in later[:i_index]:
                rows[max(i,j)].add(min(i,j))  # fill-in from elimination
    return [sorted(row) for row in rows]
\end{lstlisting}

\paragraph{What this does.}
Start with the lower-triangular part of $\mathbf{H}$: there is a block at
$(i, j)$ (with $i > j$) for every joint connecting bodies $i$ and $j$.
Then simulate Cholesky elimination: when body $k$ is eliminated, any two
bodies $i > k$ and $j > k$ that both had a block in column $k$ become
coupled, so a fill-in block $(i, j)$ is created.  For a simple chain (tree)
with bodies numbered root-to-leaf, Cholesky produces no fill-in.  For a loop
(e.g.\ a ring), there may be fill-in.

The resulting data structures (stored as Warp GPU arrays) are:

\begin{center}
\begin{tabular}{lp{8.5cm}}
\toprule
Array & Meaning \\
\midrule
\texttt{articulation\_body\_offsets} & Start/end index in the body list for each articulation. \\
\texttt{articulation\_joint\_offsets} & Start/end index in the joint list for each articulation. \\
\texttt{articulation\_bodies} & Global body IDs, grouped by articulation. \\
\texttt{articulation\_joints} & Global joint IDs, grouped by articulation. \\
\texttt{articulation\_block\_row\_offsets} & CSR-style row pointers into the block value array. \\
\texttt{articulation\_block\_cols} & Column indices for each non-zero block. \\
\texttt{articulation\_diag\_slots} & Index of the diagonal block for each body-row. \\
\texttt{body\_articulation\_local} & Local (within-articulation) index of each body. \\
\bottomrule
\end{tabular}
\end{center}

Three working arrays are also allocated (zeroed each iteration):

\begin{center}
\begin{tabular}{ll}
\toprule
Array & Purpose \\
\midrule
\texttt{values} & Block matrix entries $[\mathbf{H}_{ij}]$, one \texttt{mat66f} per non-zero block. \\
\texttt{rhs}    & Right-hand-side $\mathbf{g}_i$, one \texttt{vec6f} per body. \\
\texttt{delta}  & Solution $\boldsymbol{\delta}_i$, one \texttt{vec6f} per body. \\
\bottomrule
\end{tabular}
\end{center}

% =====================================================================
\subsection{Step 2: accumulate contact terms (body-diagonal only)}
\label{sec:contacts}

\textbf{File}: \texttt{solver\_vbd.py}, method
\texttt{\_solve\_rigid\_body\_iteration\_block\_sparse\_joints}

Before assembling the coupled joint system, the solver accumulates contact
forces and Hessians into per-body arrays using the existing contact kernels
(\texttt{accumulate\_body\_body\_contacts\_per\_body},
\texttt{accumulate\_body\_particle\_contacts\_per\_body}).  These write into:
\begin{itemize}[nosep]
  \item \texttt{body\_forces} and \texttt{body\_torques} — the contact gradient.
  \item \texttt{body\_hessian\_ll}, \texttt{body\_hessian\_al},
        \texttt{body\_hessian\_aa} — the $3\times 3$ sub-blocks of the contact
        Hessian ($\mathbf{H}_\text{lin-lin}$, $\mathbf{H}_\text{ang-lin}$,
        $\mathbf{H}_\text{ang-ang}$).
\end{itemize}

These are then merged into the \emph{diagonal} blocks of the block-sparse
matrix.  Off-diagonal contact coupling (between two bodies in contact) is
intentionally excluded — it would require extending the sparsity pattern
dynamically with every collision event, which is expensive and unnecessary for
typical contact scenarios.

% =====================================================================
\subsection{Step 3: assemble the block-sparse system (per-articulation kernel)}
\label{sec:assemble}

\textbf{File}: \texttt{rigid\_sparse\_articulation\_kernels.py},
kernel \texttt{solve\_articulation\_sparse\_cpu}

This is the main Warp kernel.  It is launched with one thread per
articulation, so each articulation is processed independently (and therefore
in parallel).  Each thread does the following.

\subsubsection{3a: inertia diagonal blocks}

For each body $i$ in the articulation, the kernel adds the inertia
contribution to its diagonal block $\mathbf{H}_{ii}$:

\begin{equation}
  \mathbf{H}_{ii}^\text{inertia}
  = \begin{pmatrix}
      m/\Delta t^2 \, \mathbf{I}_3 + \mathbf{H}^\text{contact}_{ll}
        & (\mathbf{H}^\text{contact}_{al})^\top \\
      \mathbf{H}^\text{contact}_{al}
        & \mathbf{R}\mathbf{I}_B\mathbf{R}^\top/\Delta t^2 + \mathbf{H}^\text{contact}_{aa}
    \end{pmatrix}
\end{equation}

where $m$ is the body mass, $\mathbf{I}_B$ is the inertia tensor in the body
frame, $\mathbf{R}$ is the current rotation matrix, and $\Delta t$ is the time
step.  A small regulariser ($10^{-6}$) is added to the angular diagonal to
prevent singular matrices for bodies with near-zero inertia.

The right-hand side $\mathbf{g}_i$ is initialised with the inertia force and
torque:
\begin{align}
  \mathbf{f}_i &= \frac{m}{\Delta t^2}(\mathbf{p}^\star_i - \mathbf{p}_i)
                  + \mathbf{f}^\text{contact}_i, \\
  \boldsymbol{\tau}_i &= \frac{1}{\Delta t^2}
    \mathbf{R}\,\mathbf{I}_B\,\boldsymbol{\theta}
                  + \boldsymbol{\tau}^\text{contact}_i,
\end{align}
where $\mathbf{p}^\star$ is the inertia-integrated position and $\boldsymbol{\theta}$
is the rotation-vector from the current to the inertia rotation.

\subsubsection{3b: joint contributions}

For each joint in the articulation the kernel assembles both the force/torque
gradient and the coupled Hessian blocks.

\paragraph{Linear (translational) constraint.}
Let $\mathbf{a}_p$ and $\mathbf{a}_c$ be the world-space anchor points on the
parent and child bodies.  The constraint violation is $\mathbf{C}_\text{lin} =
\mathbf{a}_c - \mathbf{a}_p$.  After projecting out the free directions with
projector $\mathbf{P}$ (identity for fixed joints, complement of the free axis
for revolute/prismatic joints), the force contributed to the RHS is
\begin{equation}
  \mathbf{f}_\text{lin} = k_\text{lin} \,\mathbf{P}\,
    (\mathbf{C}_\text{lin} - \alpha \mathbf{C}^0_\text{lin})
    + k_{d,\text{lin}} \,\mathbf{P}\, \dot{\mathbf{C}}_\text{lin}
    + \mathbf{P}\,\boldsymbol{\lambda}_\text{lin},
  \label{eq:f_lin}
\end{equation}
where $k_\text{lin}$ is the penalty stiffness, $k_{d,\text{lin}}$ is damping,
$\boldsymbol{\lambda}_\text{lin}$ is the Augmented Lagrangian multiplier, and
$\alpha, \mathbf{C}^0$ are stabilisation terms (AVBD warmstart).

The contribution to the Hessian is computed via the Jacobian $\mathbf{J}$ of
$\mathbf{C}$ with respect to the $6$-DOF body velocity.  For a body with COM
at $\mathbf{c}$ and anchor at $\mathbf{a}$, the Jacobian row maps body DOF
$d$ (linear: $d\in\{0,1,2\}$, angular: $d\in\{3,4,5\}$) to:
\begin{align}
  J_{d<3} &= \pm P_{c,d}, \\
  J_{d\geq 3} &= \pm \sum_k P_{c,k} \, [\mathbf{r}\times]_{k,d-3},
\end{align}
where $\mathbf{r} = \mathbf{a} - \mathbf{c}$ is the moment arm and $[\mathbf{r}\times]$
is the skew-symmetric matrix of $\mathbf{r}$.  The Hessian block contribution
for a pair $(i,j)$ is $\mathbf{H}_{ij} \mathrel{+}= (k_\text{lin} + k_{d,\text{lin}}/\Delta t)\,
\mathbf{J}_i^\top \mathbf{J}_j$.

\paragraph{Angular (rotational) constraint.}
Let $\boldsymbol{\kappa}$ be the \emph{rotation vector} that aligns the child anchor
frame to the parent anchor frame (relative to rest):
\begin{equation}
  \boldsymbol{\kappa} = \text{axis-angle}\bigl(q_{\text{parent}}^{-1}
    q_{\text{child}} \cdot q_\text{rest}^{-1}\bigr).
\end{equation}
The projected angular violation is $\mathbf{P}_\text{ang}\,\boldsymbol{\kappa}$.
The angular Jacobian is $\mathbf{J}_\text{ang} \in \mathbb{R}^{3\times 3}$,
mapping body angular velocity to $\dot{\boldsymbol{\kappa}}$.  The torque and
Hessian are assembled analogously to the linear case.

\paragraph{Drive and limit terms.}
For actuated joints (revolute, prismatic, D6) the kernel also assembles the
contribution of PD drives and joint limits.  For a revolute joint with axis
$\hat{\mathbf{a}}$:
\begin{equation}
  \theta = \hat{\mathbf{a}}^\top \boldsymbol{\kappa} + \theta_\text{rest},
  \quad
  \dot{\theta} = \hat{\mathbf{a}}^\top \dot{\boldsymbol{\kappa}},
\end{equation}
\begin{equation}
  f_\text{drive} = k_\text{drive}(\theta - \theta_\text{target})
    + k_{d,\text{drive}}(\dot{\theta} - \dot{\theta}_\text{target}).
\end{equation}
This is assembled as a rank-1 contribution to the angular Hessian blocks.

% =====================================================================
\subsection{Step 4: block Cholesky factorisation}
\label{sec:cholesky}

After all contributions are assembled, the kernel performs an in-place
block Cholesky factorisation of $\mathbf{H}$ following the standard left-looking
algorithm.  For column $k$ (body $k$):

\begin{enumerate}[nosep]
  \item Update the diagonal block:
        $\mathbf{A}_{kk} \leftarrow \mathbf{H}_{kk}
          - \sum_{s<k} \mathbf{L}_{ks} \mathbf{L}_{ks}^\top$.
  \item Compute the diagonal factor:
        $\mathbf{L}_{kk} = \mathrm{chol6}(\mathbf{A}_{kk})$,
        a dense $6\times 6$ Cholesky factorisation.
  \item For each $i > k$ with a non-zero block at $(i,k)$:
        $\mathbf{L}_{ik} = (\mathbf{A}_{ik}
          - \sum_{s<k} \mathbf{L}_{is}\mathbf{L}_{ks}^\top)
          \mathbf{L}_{kk}^{-\top}$.
\end{enumerate}

The $6\times 6$ dense Cholesky (\texttt{\_cholesky66}) uses a standard
lower-triangular Cholesky with a small clamp ($10^{-9}$) on the diagonal to
guard against near-singular blocks.

Block indices that are zero (no joint) are skipped using
\texttt{\_find\_block\_slot}, which searches the sparse column list.

% =====================================================================
\subsection{Step 5: block triangular solves (forward and back substitution)}
\label{sec:solve}

With $\mathbf{H} = \mathbf{L}\mathbf{L}^\top$, the Newton step is solved as:

\begin{enumerate}[nosep]
  \item \textbf{Forward substitution}: solve $\mathbf{L}\,\mathbf{y} = \mathbf{g}$
        for $\mathbf{y}$, row by row:
        \[
          \mathbf{y}_i = \mathbf{L}_{ii}^{-1}
            \!\left(\mathbf{g}_i - \sum_{j<i} \mathbf{L}_{ij}\,\mathbf{y}_j\right).
        \]
  \item \textbf{Backward substitution}: solve $\mathbf{L}^\top\,\boldsymbol{\delta} = \mathbf{y}$
        for $\boldsymbol{\delta}$, reverse order:
        \[
          \boldsymbol{\delta}_i = \mathbf{L}_{ii}^{-\top}
            \!\left(\mathbf{y}_i - \sum_{j>i} \mathbf{L}_{ji}^\top\,\boldsymbol{\delta}_j\right).
        \]
\end{enumerate}

Each $\mathbf{L}_{ii}$ solve is done with \texttt{\_solve\_lower66} /
\texttt{\_solve\_upper66\_from\_lower}, simple $6\times 6$ triangular solves.

% =====================================================================
\subsection{Step 6: apply the position update}
\label{sec:update}

The solved $\boldsymbol{\delta}_i = (\Delta\mathbf{x}_i,\,
\boldsymbol{\omega}_i) \in \mathbb{R}^6$ is applied to body $i$'s pose with
under-relaxation $\alpha_r \in (0, 1]$ (default $0.65$):

\begin{align}
  \Delta\mathbf{x}_i^\text{relax} &= \alpha_r \,\Delta\mathbf{x}_i, \\
  \mathbf{p}_i^\text{new} &= \mathbf{p}_i + \Delta\mathbf{x}_i^\text{relax}, \\
  \mathbf{r}_i^\text{new} &= \exp(\alpha_r\,\boldsymbol{\omega}_i) \cdot \mathbf{r}_i,
\end{align}
where $\exp(\boldsymbol{\omega})$ converts a rotation vector to a quaternion.

The relaxation $\alpha_r < 1$ prevents overshoot when the contact Hessian is
stale (contacts are only re-detected every few substeps), which can make the
combined inertia+contact Hessian a poor approximation.

% =====================================================================
\subsection{Step 7: update dual variables (AVBD multipliers and contact state)}

After the position update, the standard AVBD dual-variable update kernels
are called to update:
\begin{itemize}[nosep]
  \item $\boldsymbol{\lambda}_\text{lin}$, $\boldsymbol{\lambda}_\text{ang}$ —
        the Augmented Lagrangian multipliers for joints.
  \item Contact sticking flags, friction multipliers.
\end{itemize}
This is identical to the local solve path; the block-sparse and local solves
differ only in Steps 3--6.

% =====================================================================
\section{Data flow summary}
% =====================================================================

\begin{center}
\begin{tabular}{lp{10cm}}
\toprule
Phase & What happens \\
\midrule
Construction & Symbolic Cholesky pattern computed; sparse index arrays allocated. \\
Per substep & Contact kernels accumulate body-diagonal forces/Hessians. \\
Per VBD iteration & \texttt{solve\_articulation\_sparse\_cpu} kernel (one thread per articulation): \\
 & \quad 1.\ Zero sparse buffers. \\
 & \quad 2.\ Fill diagonal blocks (inertia + contact Hessian). \\
 & \quad 3.\ Fill diagonal + off-diagonal blocks (joint terms). \\
 & \quad 4.\ Block Cholesky factorisation (in place). \\
 & \quad 5.\ Forward + backward substitution. \\
 & \quad 6.\ Apply relaxed pose update. \\
Post iteration & Dual updates ($\boldsymbol{\lambda}$, contact state). \\
\bottomrule
\end{tabular}
\end{center}

% =====================================================================
\section{Contacts in the coupled system}
% =====================================================================
\label{sec:contacts_detail}

Contacts between two bodies (body-body) or between a body and a particle
(body-particle) are included via their \emph{diagonal} Hessian and force
contributions only.  They are assembled by the existing
\texttt{accumulate\_body\_body\_contacts\_per\_body} kernel before the joint
kernel runs, and the results ($\mathbf{H}^\text{contact}_{ll}$,
$\mathbf{H}^\text{contact}_{al}$, $\mathbf{H}^\text{contact}_{aa}$) are merged
into the diagonal $6\times 6$ block of each body.

This means contacts effectively couple with all joints in the articulation
(via the off-diagonal blocks), but two bodies in contact with each other do
not get an off-diagonal block between them from the contact.  This is a
deliberate design choice: the contact graph changes every substep and adding
it to the sparsity pattern would require re-running the symbolic Cholesky.
Results show this is sufficient for humanoid contact scenarios.

% =====================================================================
\section{Supported joint types}
% =====================================================================

The sparse kernel handles the following joints:

\begin{center}
\begin{tabular}{lp{4.5cm}p{4.5cm}}
\toprule
Joint & Linear part & Angular part \\
\midrule
Fixed      & Full 3D penalty & Full 3D penalty \\
Revolute   & Full 3D penalty & Penalty on 2 perp.\ axes; drive on free axis \\
Prismatic  & Penalty on 2 perp.\ axes; drive on free axis & Full 3D penalty \\
Ball       & Full 3D penalty & (no angular penalty) \\
D6         & Per-axis mix & Per-axis mix \\
Cable      & One-sided tension along cable axis & Bend penalty with Dahl friction \\
\bottomrule
\end{tabular}
\end{center}

\emph{Free-direction projection} is done with the projector
$\mathbf{P} = \mathbf{I} - \hat{\mathbf{a}}\hat{\mathbf{a}}^\top$ (identity
minus the outer product of the free axis unit vector).  This removes the free
direction from the constraint force, so only the constrained DOFs are penalised.

% =====================================================================
\section{Performance and convergence}
% =====================================================================

\subsection{Residual reduction}

On 128-body synthetic chains (fixed or revolute joints) and loops, the sparse
articulation solve achieves:
\begin{itemize}[nosep]
  \item $1.3$--$1.6\times$ one-iteration residual reduction vs the local solve.
  \item $\sim 2\times$ residual reduction after 8 VBD iterations.
  \item $117\times$ energy reduction on an 8-body fixed ring with a
        $10^6$-to-$1$ stiffness ratio (linear vs.\ angular stiffness).
\end{itemize}

Prismatic joints benefit the most from coupling (16--18$\times$ linear energy
reduction); D6 joints show more modest gains ($2.6\times$ linear).

\subsection{Timing (CPU, 44-body Unitree G1)}

\begin{center}
\begin{tabular}{ll}
\toprule
Solver mode & Mean time per iteration \\
\midrule
Local (diagonal) solve & 51.2 $\mu$s \\
Block-sparse solve (serial CPU) & 73.1 $\mu$s \\
\bottomrule
\end{tabular}
\end{center}

The coupled solve is $\sim 1.4\times$ slower per iteration on CPU.  Because it
converges faster, the net cost per unit of residual reduction is lower.

A CUDA block-sparse kernel (GPU) is noted as future work; the current
implementation only runs on CPU.

% =====================================================================
\section{How to use the new feature}
% =====================================================================

Enable the sparse articulation solve by passing the argument
\begin{center}\texttt{rigid\_articulation\_solve="block\_sparse\_joints"}\end{center}
to \texttt{SolverVBD}:

\begin{lstlisting}
solver = SolverVBD(
    model,
    rigid_articulation_solve="block_sparse_joints",
    rigid_articulation_relaxation=0.65,   # default, tune if needed
    sim_substeps=10,
    sim_iterations=5,
)
\end{lstlisting}

\paragraph{Notes.}
\begin{itemize}[nosep]
  \item Only available on CPU for now (\texttt{device="cpu"}).
  \item \texttt{rigid\_articulation\_relaxation} $\in (0, 1]$. Lower values
        ($\sim 0.5$) improve stability when contacts are stale; higher values
        ($1.0$) give faster convergence when contacts are well-resolved.
  \item The baseline \texttt{"local"} mode remains the default.
\end{itemize}

% =====================================================================
\section{Limitations and future work}
% =====================================================================

\begin{itemize}
  \item \textbf{CPU only.}  CUDA block-sparse kernels are not yet implemented.
        The serial CPU kernel runs one articulation per thread; a CUDA version
        would need warp-level parallelism over bodies within an articulation.
  \item \textbf{No off-diagonal contact coupling.}  Contacts between bodies of
        the same articulation only appear in diagonal blocks.  In scenarios with
        self-contact (robot hitting itself) this may limit convergence.
  \item \textbf{Static sparsity pattern.}  The symbolic Cholesky is computed
        once at construction.  Adding/removing joints at runtime would require
        rebuilding the layout.
  \item \textbf{May require relaxation for nonsmooth contact.}  Stiff contact
        with poor contact Hessian approximation can cause divergence at
        relaxation $= 1.0$.
\end{itemize}

% =====================================================================
\section{Quick-reference: key files}
% =====================================================================

\begin{center}
\begin{tabular}{p{7.5cm}p{6cm}}
\toprule
File & Responsibility \\
\midrule
\texttt{rigid\_sparse\_articulation.py} &
  Sparse layout construction (once at init). \\
\texttt{rigid\_sparse\_articulation\_kernels.py} &
  Warp kernel: assembly, Cholesky, solve, update. \\
\texttt{solver\_vbd.py} &
  Dispatches local vs.\ block-sparse solve. \\
\texttt{tests/test\_vbd\_sparse\_articulation.py} &
  Tests: chain, loop, contacts, drives. \\
\texttt{reports/vbd\_sparse\_articulation/} &
  Benchmarks and validation scripts. \\
\bottomrule
\end{tabular}
\end{center}

\end{document}
